\section{Personlige Refleksioner}
\subsection{Jonathan}
% jonathan
Jeg tog en mindre rolle i projektets opstart og dette påvirkede mit engagement i de tidlige faser af projektet, og jeg oplevede også perioder, hvor jeg ikke deltog i møder eller informerede gruppen om fravær på forhånd. Set i bakspejlet kunne jeg have bidraget mere aktivt. Jeg tror, at min manglende deltagelse i starten kan have været en form for social loafing, hvor jeg ubevidst trak mig tilbage fra ansvaret, fordi jeg følte, at andre i gruppen ville tage sig af opgaverne, eller allerede havde taget sig af dem. For at bryde ud af denne dynamik forsøgte jeg at bidrage på områder, hvor jeg oplevede, at der manglede arbejdskraft, blandt andet gennem kildearbejde og rapportskrivning. Da jeg som person har en tendens til at være tilbageholdene i sociale sammenhænge, var dette en måde jeg stadig kunne bidrage på.
Jeg oplevede at gruppen håndterede situationen professionelt og inkluderende, hvilket gjorde det lettere for mig at bidrage mere aktivt senere i projektforløbet. Jeg har lært, at det er vigtigt at kommunikere åbent om ens deltagelse og forpligtelser i et gruppeprojekt for at sikre, at alle medlemmer føler sig ansvarlige og engagerede. I fremtidige projekter vil jeg stræbe efter at være mere proaktiv i min deltagelse og kommunikere klart med mine gruppemedlemmer om mine bidrag og eventuelle udfordringer, jeg måtte have.
\newpage

\subsection{Victor}
I forhold til selve implementeringsfasen af projektet endte jeg med at tage en meget specifik vinkel på projektet, nemlig implementeringen af voicechat-systemet. Dette endte dog med at være det eneste, jeg bidrog med i forhold til implementering, grundet uforudsete komplikationer. Set i bakspejlet burde jeg have været bedre til at bede om hjælp til visse aspekter af arbejdet, således at jeg kunne have arbejdet på flere dele af spillet frem for kun én enkelt.

Jeg har til tider været sent ude med mine svar i gruppen og missede også et enkelt scrum-møde, hvilket er beklageligt. Især mod slutningen af projektet låste jeg mig fast på at få voicechat-systemet til at fungere, hvilket førte til forsinkede GitHub-pushes, da jeg havde et stort fokus på, at det system, jeg arbejdede på, skulle være fuldt funktionelt.

Den vigtigste læring, jeg tager med fra dette projekt, er især ikke at hyperfiksere på et bestemt problem, men i stedet at blive bedre til at sprøge om hjælp til- og fordele større arbejdsopgaver blandt gruppemedlemmer. På den måde kan jeg fremover få større indflydelse på flere dele af udviklingen og dermed også opnå et bedre indblik i projektets forskellige systemer.
\newpage

\subsection{Nikolaj}
Som Scrum Master har jeg fra starten af projektet påtaget mig et stort ansvar i gruppen. Dette har passet godt til min personlige arbejdsstil, som er præget af en høj grad af ansvarstagen og et ønske om at skabe fremdrift hurtigt i projektet. Ud fra et personligt perspektiv (jf. DISC-profil, rød type) trives jeg i roller, hvor jeg kan tage lederskab og sikre progression i arbejdsprocessen.

I løbet af projektforløbet har jeg dog ofte oplevet, at en stor del af fremdriften og koordineringen i praksis har været båret af mig selv. Flere gruppemedlemmer har i de tidlige faser haft udfordringer med deres opgaver, hvilket har medført, at jeg i mange tilfælde har måttet overtage, støtte eller færdiggøre dele af arbejdet for at sikre, at projektet kunne fortsætte. Dette har til tider haft en demotiverende effekt, da jeg har følt, at jeg i høj grad stod alene med ansvaret for projektets progression.

Samtidig må jeg også erkende, at jeg selv kunne have håndteret situationen anderledes. I stedet for primært at “løbe foran” og løse opgaverne selv, kunne jeg have været mere aktiv i forhold til at facilitere samarbejde og sikre, at gruppemedlemmerne blev hjulpet videre i deres opgaver. Dette kunne eksempelvis være gjort gennem tydeligere opfølgning, bedre fordeling af ansvar samt mere aktiv involvering af gruppen i problemløsning. Set i et PBL-perspektiv kan dette relateres til vigtigheden af fælles ansvar for læring og vidensdeling i gruppen.

Arbejdssituationen har i perioder været præget af høj belastning, hvor jeg har arbejdet mange timer og haft begrænset hvile for at sikre projektets fremdrift. Dette har været energimæssigt krævende og til tider udmattende. En del af udfordringen skyldes, at vi relativt sent i forløbet kom i gang med egentlig implementering, hvilket skabte et presset udviklingsforløb senere i projektet.

Samtidig oplevede jeg, at opgavefordelingen i gruppen til tider var vanskelig, da nogle medlemmer holdt sig tilbage i forhold til at påtage sig bestemte opgaver, hvilket muligvis skyldtes usikkerhed omkring egne kompetencer. Dette førte til en skæv arbejdsfordeling, hvor jeg i høj grad endte med at varetage både udvikling, strukturering af projektet, rapportskrivning, mødeplanlægning samt håndtering af dokumentation og kilder.

Jeg har også erfaret, at kommunikationen i gruppen i nogle tilfælde var afhængig af formelle Scrum-møder, og at der manglede proaktiv opfølgning mellem disse møder. Dette har gjort, at problemer i perioder ikke blev håndteret tidligt nok, hvilket påvirkede effektiviteten i arbejdet.

Et vigtigt læringspunkt er, at jeg fremadrettet skal blive bedre til at balancere mellem at tage ansvar og samtidig sikre, at ansvar faktisk bliver fordelt og fastholdt i gruppen. Dette indebærer blandt andet en mere konsekvent brug af opfølgning, tydeligere forventningsafstemning og en større grad af facilitering frem for direkte overtagelse af opgaver.

Derudover har projektet givet mig indsigt i, hvor vigtigt det er at starte implementering tidligere i processen samt hvor afgørende det er at kunne tilpasse sig eksterne afhængigheder, såsom samarbejdet med Viborg Katedralskole, der også delvist er skyld i den sene implementation start, da der var planlagt og strukktueret et markant større sammenarbejde med Skolen, dog vist de sig at trække meget på aftale og planer, og derfor desværre ikke forløb som forventet og derfor ikke i fuldt omfang blev integreret i det endelige produkt.

Samlet set har projektet givet mig betydelig erfaring med både teknisk projektstyring, samarbejde i en PBL-kontekst og personlig ledelse. Fremadrettet vil jeg arbejde på at forbedre min evne til at delegere, støtte gruppens selvstændige læring og skabe en mere balanceret arbejdsfordeling, så jeg ikke alene bærer hoveddelen af projektets fremdrift.